%% =====================================================================
%%  B5-T-14-01 -- what B5 contributes to "The Automatic Response"
%%  -------------------------------------------------------------------
%%  LAYER A: APPEND-ONLY. Written once, then left alone. Material found
%%  only here sits in the `unique` slot and is protected by rule R10.
%% =====================================================================
%% @kind: source
%% @id: B5-T-14-01
%% @book: B5
%% @topic: T-14-01
%% @locator: ch. 1, pp. 13--35
%% @status: distilled
%% @unique: yes
%% @integrated: yes
%% @updated: 2026-09-19

\dossier{B5}{T-14-01}{\bookshort{B5}}{ch. 1, pp. 13--35}

\begin{distilled}
  \item The framing chapter restates negotiation as life's operating system and grounds the whole method in the two-systems picture: a fast, emotional system that reacts first and a slow, rational one that arrives late --- under pressure, the fast system answers.
  \item Its named method for dealing with that architecture is tactical empathy plus calibrated questions: calm the emotional system rather than out-argue it.
  \item The chapter's evidence base is announced as hostage negotiation plus the behavioural-economics literature, positioning the whole book as applied automatic-response engineering.
  \item It argues the assumption that counterparts are rational bargainers is the amateur error; they are reactive systems that become rational only after the emotional load is handled.
\end{distilled}

\begin{unique}
  \item The two-systems premise stated as the practitioner's first diagnostic: which system is answering right now, theirs or yours --- and the technique set is chosen by the answer.
  \item The de-escalation-first doctrine: since the fast system cannot be argued with, every approach begins by lowering its arousal before any content is exchanged.
\end{unique}
